\section{Pricers}\label{sec:arch-pricers-chapter}
A pricer turns an instrument and the \code{MarketEnvironment} into a float64 price, as one TensorFlow graph: the greeks
are its gradients with respect to the market variables. One subsection for the common interface, then one per asset
class.

% ─────────────────────────────────────────────────────────────────────────────
\subsection{Pricer interface and factory}\label{sec:arch-pricers}
\code{Pricer}: \code{calculate\_price}, \code{price(..., autodiff)}, the tape; \code{PricerAssignment}.

\paragraph{\code{Pricer} (\code{pricers/pricer.py}).} A pricer implements \code{calculate\_price(product,
market\_env)}: a TensorFlow graph from the market objects to a float64 scalar. \code{price(product, market\_env,
autodiff=False)} runs it inside \code{market\_env.pinned()} (the monotone convex coefficients of the curves computed
once per call, inside the tape, so the gradients are unchanged), stores the price on \code{product.price} and, with
\code{autodiff=True}, keeps the \code{GradientTape} on \code{pricer.tape}:
\begin{itemize}[nosep]
  \item greeks: \code{pricer.tape.gradient(product.price, variables)} for any market variables, or for all of them
        with \code{market\_env.risk\_factors()} (bucketed rho per pillar, vega per node, delta, dividend
        sensitivities in one call);
  \item second orders (gamma): a second tape around the call;
  \item model greeks: the gradient with respect to \code{model.parameters} when the pricer is given a model instance
        (Section~\ref{sec:arch-equity-models}).
\end{itemize}
Rules: the market data are read only through the \code{MarketEnvironment} accessors (or the \code{Equity} view), never
copied into new variables nor detached with \code{.numpy()} / \code{float()} (the gradient would stop there); the only
variable a pricer may create is the time to maturity, left on the product (\code{product.time\_to\_maturity}) for
theta. The pricers leave their intermediate values on the product for inspection (forward, vol, discount factor,
annuity, standard error, \dots); they are diagnostics, not inputs.

\paragraph{\code{PricerAssignment} (\code{pricers/factory.py}).} \code{PricerAssignment.create(product)} returns a new
default pricer by \code{isinstance} (subclasses covered), \code{TypeError} for an unknown product:
\begin{center}
\small
\begin{tabular}{@{}L{5.4cm}L{9.8cm}@{}}
\toprule
Product & Default pricer \\
\midrule
\code{Deposit}, \code{Fra}, \code{Swap}, \code{Ois} & \code{DepositPricer}, \code{FraPricer}, \code{SwapPricer},
  \code{OisPricer} \\
\code{CapFloor}, \code{Collar} & \code{CapFloorPricer} \\
\code{Swaption} & \code{SwaptionPricer} \\
\code{VanillaOption}, \code{DigitalOption}, \code{BarrierOption} & \code{BlackScholesPricer()} (default arguments) \\
\code{AutocallableOption} & none: \code{AutocallableMCPricer} chosen explicitly \\
\bottomrule
\end{tabular}
\end{center}
The Monte Carlo pricers are always chosen explicitly, with their model and simulation settings.

\paragraph{Adding a pricer.} Subclass \code{Pricer}, implement \code{calculate\_price} on the accessors, raise
\code{TypeError} for a product it does not price and \code{ValueError} for missing market data; register it in
\code{PricerAssignment} if it is the default for its product; test it against an independent check (closed form,
parity, Monte Carlo within a few standard errors) and its AAD greeks against bumps.

% ─────────────────────────────────────────────────────────────────────────────
\subsection{Interest-rate pricers}\label{sec:arch-rates-pricers}
Discounting pricers for the linear products and closed forms on the rate vols for the options.

\paragraph{Curves.} Every rate pricer discounts on the default curve of the product's currency
(\code{market\_env.get\_ir\_curve(product.ccy)}, the overnight one: ESTR, SOFR, SONIA) and projects an Ibor index on the
curve named by the index (\code{EUR:6M} $\to$ \code{get\_ir\_curve(EUR, ticker="6M")}); an OIS uses the overnight
curve for both. The curve names used are left on the product (\code{discount\_curve}, \code{estimation\_curve}). A
missing curve raises \code{ValueError}.

\paragraph{Linear products.}
\begin{itemize}[nosep]
  \item \code{DepositPricer} (\code{deposit.py}): $+N\,P(s)$ at the start (if not past) $-\,N(1 + r\tau)\,P(e)$ at the
        end;
  \item \code{FraPricer} (\code{fradiscounting.py}): settled at the start,
        $\text{side}\cdot N\tau(F - K)\,P_{\text{disc}}(s)/(1 + F\tau)$, $F$ on the index curve; zero once started;
  \item \code{SwapPricer}, \code{OisPricer} (\code{swapdiscounting.py}): the two legs priced by the leg discounting
        classes, signed by \code{swap\_type} (payer: floating $-$ fixed); the leg PVs are left on
        \code{fixed\_leg.price} and \code{floating\_leg.price};
  \item leg discounting (\code{fixedflow.py}, \code{floatingflow.py}): \code{FixedCouponDiscounting} /
        \code{FixedLegDiscounting} (amounts $\times$ discount factors, vectorised with \code{PrefetchedCurve}),
        \code{FloatingCouponDiscounting} / \code{FloatingLegDiscounting} (Ibor forward on the estimation curve,
        discounted on the discount curve), \code{OisCouponDiscounting} / \code{OisLegDiscounting} (compounded
        overnight rate, single curve); a fixing already past is read from the index, a paid coupon is worth zero.
\end{itemize}
The curve bootstrap prices its instruments with these same pricers.

\paragraph{\code{CapFloorPricer} (\code{capfloor.py}).} Caplet by caplet, $N\tau\,P_{\text{disc}}(\text{pay})\times
\text{option}(F, K, \sigma, T_{\text{fix}})$: forward on the index curve, vol from the caplet surface
\code{get\_ir\_vol\_surface(ccy, tenor)} at (fixing time, strike), \code{bachelier\_calc} for \code{vol\_type}
\code{"normal"}, \code{black\_calc} for \code{"lognormal"}; a caplet already fixed pays its intrinsic value on the
index fixing, a paid one is dropped. A \code{Collar} is cap $-$ floor (prices also on \code{collar.cap.price} and
\code{collar.floor.price}); the caplet values on \code{product.caplet\_values}.

\paragraph{\code{SwaptionPricer} (\code{swaption.py}).} $A\times\text{option}(F, K, \sigma, T)$, annuity $A$ and
forward swap rate $F$ from \code{swap\_annuity\_and\_forward} (two curves, as \code{SwapPricer}), vol from the cube
\code{get\_swaption\_vol\_cube(ccy, tenor)} at (expiry, \code{tenor\_years}, $K - F$), Bachelier or Black-76 by the
cube's \code{vol\_type}. Diagnostics: \code{annuity}, \code{forward\_swap\_rate}, \code{volatility}.

The rate Monte Carlo (\code{GaussianShortRateModel}) is not a \code{Pricer}: its closed forms (bonds, caplets,
swaptions) serve the calibration and the exposures (Section~\ref{sec:arch-models-chapter},
\code{SwapExposureGenerator} in \code{analytics/}).

% ─────────────────────────────────────────────────────────────────────────────
\subsection{Equity pricers}\label{sec:arch-equity-pricers}
Every equity pricer reads the underlying through \code{market\_env.get\_equity(product.underlying)}: spot, forward,
dividends, repo, vol surface and discount curve from one place (Section~\ref{sec:arch-equity}).

\paragraph{\code{BlackScholesPricer} (\code{black.py}).} Closed forms, as functions of float64 tensors:
\code{blackscholes\_calc} (European), \code{bjerksund\_stensland\_calc} (American, 1993), \code{digital\_calc} (cash- /
asset-or-nothing), \code{barrier\_calc} (Reiner--Rubinstein, knock-in at expiry exact, Broadie--Glasserman--Kou shift
for discrete monitoring). \code{BlackScholesPricer(dividend\_model="continuous", use\_implied\_repo=True,
digital\_spread=None)} builds the inputs from the \code{Equity} view (\code{\_build\_inputs}: spot, repo to the
expiry, rate and discount factor of the currency curve, vol of the surface at $(K, T)$, dividends) and dispatches on
the product: \code{BarrierOption} $\to$ \code{barrier\_calc}, \code{DigitalOption} $\to$ \code{digital\_calc} or, with
\code{digital\_spread} $= h$, a centred spread of vanillas at $K(1 \pm h)$ each at its own vol (the smile slope
included), \code{VanillaOption} $\to$ European or American. \nb{black\_scholes.ipynb}, \nb{digital\_options.ipynb}

\paragraph{Dividends and carry}\label{sec:arch-bs-dividends} With a discrete dividend schedule on the underlying,
\code{dividend\_model="discrete"} prices on \code{spot\_net} $= S - $ \code{pv\_dividends(T)} with carry = repo;
\code{"continuous"} turns the schedule into $q_{\mathrm{eff}}$, computed with the spot under \code{stop\_gradient}, so
the AAD delta is the sticky-yield one (a bump-and-reprice recomputes $q_{\mathrm{eff}}$ and gives the sticky-cash
delta). Without a schedule the market \code{DIVYIELD} is used whatever \code{dividend\_model}.
\code{use\_implied\_repo=False} drops the repo from the carry. The pricer leaves its inputs on the product:
\code{spot}, \code{spot\_net}, \code{pv\_discrete\_dividends}, \code{repo\_margin}, \code{risk\_free\_rate},
\code{discount\_factor}, \code{volatility}, \code{time\_to\_maturity} (a variable: theta), \code{forward}.
\nb{dividends.ipynb}

\paragraph{Monte Carlo pricers: model class or instance.} Both Monte Carlo pricers take a \code{model} argument:
\begin{itemize}[nosep]
  \item a \emph{class} with \code{from\_equity} (\code{GeometricBrownianMotion}, \code{LocalVolatilityModel},
        \code{DisplacedDiffusionModel}, \code{HestonModel}, \code{HestonSLVModel}): the process is built and
        calibrated from the market at every call, inside the tape, so the market greeks go through the calibration;
  \item an \emph{instance} (\code{EquityModel}, e.g. loaded from a calibration set): bound with
        \code{model.process(equity, maturity)}, parameters held; model greeks on \code{model.parameters}, par greeks
        with \code{par\_greeks}.
\end{itemize}
The normals are drawn by the pricer with \code{tf.random.stateless\_normal} (\code{seed}),
\code{model.noise\_size()} per path and step, and passed to \code{evolve} (Section~\ref{sec:arch-models}): the same
seed gives the same paths, so bumps are common-random-number ones. The digital and barrier indicators are smoothed by
a sigmoid (width as a fraction of the strike or barrier, 1\% by default, 0 = hard indicator), without which their
pathwise AAD greeks would be zero.

\paragraph{\code{VanillaMCPricer} (\code{vanillamc.py}).} European \code{VanillaOption}, \code{DigitalOption} and
\code{BarrierOption} (American exercise refused). Arguments: \code{model} (default GBM), \code{n\_paths} (50\,000),
\code{n\_steps} (252, uniform grid ending at the expiry), \code{seed}, \code{digital\_smoothing}. Barriers: continuous
monitoring by the Brownian-bridge survival probability on each step (vol of the surface at $(H, T)$; no smoothing
needed), discrete monitoring on the monitoring dates added to the grid. After pricing: \code{standard\_error},
\code{discount\_factor}, \code{time\_to\_maturity}.

\paragraph{Quanto (\code{quanto.py}, \code{black.py}).} \code{BlackScholesPricer}: for a quanto product the
discount rate is the payment currency's and the carry \code{q + (r\_pay - r) + rho sigma sigma\_X} (European vanilla,
digital, barrier; American rejected). Monte Carlo: \code{apply\_quanto} sets the drift on the process
(\code{StochasticProcess.set\_quanto(rho * sigma\_X)}: \code{GeometricBrownianMotion} lowers \code{mu} by it times
\code{sigma}, \code{LocalVolatilityProcess} by it times the local vol of the path and stops rescaling the paths to the
forward; the other processes raise) and returns the payment curve for the discounting; diagnostics
\code{product.quanto\_correlation}, \code{product.quanto\_fx\_vol}. Composite: \code{BlackScholesPricer} on $Y = S X$
(spot $S X_0$, forward $F_S F_X$, vol $\sqrt{\sigma_S^2 + \sigma_X^2 + 2\rho\sigma_S\sigma_X}$, payment discount);
Monte Carlo: the quanto drift and \code{composite\_fx\_paths} ($X$ lognormal on its forward, correlated with the
underlying's normals, an independent stream \code{[seed, 1]}), the paths multiplied by it (composite barrier options
rejected). \code{Pricer.price(..., price\_currency=)}: the price (and \code{product.standard\_error}) converted from
the payment currency at \code{get\_fx\_spot} inside the tape (delta to the FX spot by AAD);
\code{product.price\_currency}.

\paragraph{\code{AutocallableSwapPricer} (\code{montecarlo.py}).} Subclass of \code{AutocallableMCPricer} for an
\code{AutocallableSwap} (which \code{AutocallableMCPricer} refuses): \code{\_payoff} records the alive indicator after
each autocall payment date (\code{\_survival}), \code{\_funding} weights each funding period by it at the period start,
forward on \code{get\_ir\_curve(payment\_ccy, funding\_curve)}, discounted on \code{get\_ir\_curve(payment\_ccy)};
\code{product.funding\_pv} (per unit notional), \code{projections()} adds the \code{funding} flows. Both pricers
multiply by \code{product.side\_sign} (+1 client, -1 issuer), the projected flows too. Tradepad: the
\emph{Autocallable Swap} screen (funding index and margin), Side on both autocallable screens.

\paragraph{Vega by expiry (\code{analytics/vega.py}).} \code{vega\_by\_expiry(vega, equity)} reads a node-vega
matrix (price per vol point, \code{[n\_maturities, n\_strikes]}) by expiry: \code{level} (row sum),
\code{cumulative}, \code{skew} (the row weighted by \code{skew\_weights(equity)}: $(1 - K/F)/$\code{SKEW\_STEP}, the
moneyness clipped to \code{SKEW\_RANGE}); the Tradepad shows it next to the vega map. \nb{autocallable\_skew\_sensy.ipynb}

\paragraph{\code{AutocallableMCPricer} (\code{montecarlo.py}).} \code{AutocallableOption}, investor's side:
$N\,\E\big[\sum_i \text{coupon}_i\,P(0,\text{pay}_i) - \text{put}\,P(0,T)\big]$. Arguments: \code{model} (default
\code{LocalVolatilityModel}; \code{GeometricBrownianMotion} accepted for debugging only: the Black--Scholes model of
one expiry, wrong at the intermediate fixings), \code{n\_paths} (20\,000), \code{seed}, \code{max\_time\_step} (1/52: the grid is the
union of the future fixing dates and \code{simulation\_grid}: daily steps (\code{fine\_time\_step}) up to
\code{fine\_horizon} = 3 months, steps of \code{max\_time\_step} after; \code{fine\_horizon} = 0 gives the uniform
grid), \code{barrier\_smoothing} (1\%). Forward start: the
strike date is simulated and the initial level is $S(t_0)$ on each path. After pricing: \code{product.coupon\_pv},
\code{product.put\_pv} (per unit notional), \code{product.standard\_error}; \code{pricer.paths} at the fixings
(\code{pricer.date\_grid}, \code{pricer.time\_grid}) for diagnostics. Worst-of: \code{model} a class for every
underlying or a dict ticker $\to$ instance; one process per underlying (\code{noise\_size() = 1}: local vol, GBM)
inside a \code{MultiAssetProcess} with \code{market\_env.get\_correlation} of the \code{EQ:CCY:TICKER} names; draws
\code{[n\_paths, n\_steps, n]} (seed \code{[seed, 0]}); \code{pricer.paths} \code{[n\_paths, n\_fixings, n]},
\code{pricer.performance} the worst performance; quanto drifts on \code{product.quanto} (ticker $\to$ (correlation,
FX vol)). The payoff is written on the performance in both cases (one underlying: $S/K$). Live trades: the
observations and the strike date up to the evaluation date are known levels (\code{\_known\_levels}:
\code{product.fixings}, then \code{market\_env.get\_eq\_fixing}, the spot for a fixing on the evaluation date not in
either; a past date without a fixing raises); they enter the same payoff with exact indicators, before the simulated
dates; payments before the evaluation date are left out; with no future observation nothing is simulated
(\code{pricer.paths} \code{None}). \code{trade\_state(product, market\_env)}: strikes, per past observation the
levels, the worst performance, coupon paid / missed (with the memory) and autocall, \code{called\_on},
\code{unpaid\_coupons} (shown by the Tradepad). Composite products with past fixings are rejected.
\code{projections()}: from the paths of the last pricing (exact indicators, not discounted, per unit notional) the
expected coupons, redemptions (the notional back, at an autocall or at maturity) and puts (negative, paid by the
investor at maturity) per payment date, the
probability of the autocall at each observation and cumulated, at maturity the probability of reaching it and of the
knock-in and the loss given the knock-in, the expected life (mean time to the redemption), for a worst-of the share
of the paths where each underlying is the worst at maturity (among those reaching it and among the knocked-in ones),
and per underlying the mean simulated level at each future observation
against the market forward (a quanto underlying carries the quanto drift). The discounted coupons and puts give back the
price with the hard payoff; the Tradepad shows above the chart the price split into the PV of the coupons and of the
put (as priced). The Tradepad computes them only when its Projections tab is opened (the paths are
kept with the pricing job). \nb{digital\_strip.ipynb},
\nb{down\_and\_in\_put.ipynb}
